\chapter{Discussion}
In this chapter, we discuss the key findings within the results of the 
experiments. Then we discuss shortcomings of the research so far. Lastly,
recommendations are given. 

\section{Findings}
The results of the testing on the windowing function confirm the 
consensus in the literature, that the Hann and Hamming windows 
are both
effective in reducing spectral leakage when compared to no window. 
Through comparing the two windows, we find that the Hamming 
window is more appropriate 
for radar purposes with its sharper drop-off that allows detection to 
be more precise, at the cost of more global spectral leakage. 
Additionally, the insignificant
time difference between either window or no window is expected given
the simplicity of the function, with each kernel performing just one
multiplication. The negligible latency difference across the desktop and Xavier
platforms further indicates that window selection can be made independently of
hardware configuration. The differing spectral leakage patterns motivate a more
comprehensive search for window functions that are most appropriate for
miniature radar applications. 

The results from the coherent integration experiment find that the signal 
to noise ratio improves by the square root of the number of chirps summed 
together under Gaussian noise, creating a positive downstream effect on the 
ability of CFAR to detect weak 
targets. The difference in throughput of the processing subsystem on 
the desktop versus the Xavier is likely due to the overhead of handling data,  
such as reading and writing. The large amount of work that each thread has to 
do, summing over many chirps, motivates the exploration of more parallelised 
implementations that distribute the sums between threads. 

The CFAR results outline a method for determining a threshold from a 
desired false alarm rate. The experiments demonstrate the GPU implementation
is most appropriate for the Xavier embedded platform, regardless of the size 
of the buffer for all practical window sizes. The difference between hardware
platforms highlights the importance of hardware-aware algorithm selection. The 
results of the angle sensing stage demonstrate the capability of the algorithm 
to detect multiple targets with different bearing/elevation ratios. The ablation 
study illustrates that CFAR dominates the latency 
of the tracking subsystem. The relatively high cost of CFAR motivates the exploration
of more efficient CFAR implementations. Potential approaches include implementing 
the prefix sum table approach on the GPU, using a reduction approach to calculate the 
noise floor, and more consciously minimising cache misses within warps. 

The completed results begin to address two of the three research gaps identified in the
literature review. The GPU pipeline benchmarking across desktop and Xavier platforms, covering CFAR, angle
sensing, windowing, and coherent integration, directly addresses
gap 2, providing the hardware-validated latency, throughput, and architectural profiling
that prior UAV DAA radar work lacked. The hardware-aware implementation comparisons
partially address gap 3, establishing a baseline for angle sensing performance, though
evaluation under realistic multi-target and non-Gaussian noise conditions remains to be
completed. Gap 1, concerning detection sensitivity against small moving airborne targets,
and the remaining scope of gap 3, require the tracking stage and real flight data, which
are the subject of the remaining work.

\section{Shortcomings}
The project has progressed broadly in line with the initial proposal, with the 
preprocessing subsystem, CFAR detection, monopulse angle sensing, and the supporting 
simulation and profiling infrastructure all implemented and evaluated. The pipeline 
architecture has been validated on both synthetic data and real AWR2243 hardware data, 
and hardware-aware implementation comparisons have been completed for the CFAR stage 
across the desktop and Xavier platforms.

Several goals from the initial proposal were not completed within this reporting period.
The multi-target tracking stage, comprising clustering and the tracking filter, was not
implemented or evaluated. This represents the most significant deviation from the
proposed plan, and means that the end-to-end pipeline remains incomplete and that track
quality cannot yet be used as an evaluation metric. The primary cause of this delay was
the time required to bring up the hardware pipeline and validate each stage on real radar
data. The remaining work in the project plan is structured to address this directly.
Until the tracking stage is complete, gap 1 (detection sensitivity against small moving
airborne targets) cannot be addressed, as it requires track-level performance metrics
from real encounter scenarios.

% Incorrect below
The window function evaluation was narrowed to Hann and Hamming windows, deferring the
evaluation of Blackman, Kaiser, and Taylor windows. The coherent integration
comparison was limited to synthetic data for SNR validation and real test-source data for
hardware correctness; the parallelised implementation comparison and downstream effect on
angle sensing stability were not completed. The angle sensing evaluation did not include
the proposed sweeps over SNR, channel imbalance, and multi-target proximity. These items
are carried forward into the remaining work plan. Additionally, gap 3 — evaluating angle
sensing algorithms under realistic noise conditions — remains open until real flight data
is collected to characterise the noise regime of the AWR2243 in airborne operation. 

\section{Recommendations}
The windowing results support the selection of the Hamming window for radar applications 
where resolution of closely spaced targets is a priority, given its steeper drop-off near 
the target bin relative to the Hann window. The global spectral leakage penalty of the 
Hamming window is acceptable in practice, as the noise floor is dominated by thermal and 
clutter contributions rather than windowing artefacts at operationally relevant SNRs. As the 
latency contribution of the windowing stage is negligible and invariant to window choice, the 
selection can be made purely on detection performance grounds.

For the CFAR and angle sensing stages, the results support a GPU-based implementation on the 
Xavier embedded platform across all practical training cell window sizes. On the desktop, the 
prefix sum CPU implementation is competitive for small buffer sizes, but the GPU implementation 
is preferable at larger buffer sizes. The crossover point differs between platforms, confirming 
that implementation selection cannot be made platform-agnostically. As CFAR dominates the latency of the detection subsystem, with angle sensing contributing
a factor of 1.5 to 5 times less, optimisation effort should be directed at the CFAR stage
rather than the angle sensing stage.

The overhead of host--device data transfer and thread synchronisation constitutes a substantial 
portion of the remaining latency, and represents the primary target for further pipeline 
optimisation.
